\section{Conclusions}
\label{sec:conclusions}

\begin{enumerate}[leftmargin=1.6em,itemsep=4pt,topsep=4pt,parsep=0pt]

\item \emph{What is new.} We searched \NWinA{} fine-channel ALMA windows
toward \NSysClassA{} systems within 40\,pc for a drifting narrowband carrier at
\SurvFreqLoA--\SurvFreqHiA\,GHz: \CvUnionA\,GHz of unique frequency in
\CovNIslands{} separate intervals, at least \SbPriorRatio{} times either
earlier millimetre search of nearby stars \citep{Steffes1994,Mauersberger1996},
with \EdrNWinA{} windows long enough for the drift to cross a channel
(\S\ref{sec:sample}).

\item \emph{How deep.} The median system's best window would have registered a
carrier of $\EirpNinetySysMedA$\,W nine times in ten, and \NSysEfifteen{} of
the \NSysALim{} systems with a constrained window reach $10^{15}$\,W, or
\NmDecAdjNEfifteen{} allowing for atmospheric decorrelation; a
\BenchDiam\,m dish radiating \BenchPowerMW\,MW at \BenchFreqGHz\,GHz, beamed at
the Sun, lies above the limit for \BenchNSysA{} of them. For \SmcNEstSys{} of the \NSysALim{}
systems they are transferred from injections into other windows
(\S\ref{sec:pninetysel}). For Proxima Centauri the census stack reaches
$\SzProxEirpCensus$\,W for an emitter steady in the star's own frame;
for a transmitter on \PxPlanet{} the archive's deepest single epoch,
$\PxEirpSingle$\,W and not a census limit, applies (\S\ref{sec:stack}).

\item \emph{What was found.} Of the \EvNCrossQp{} quality-passing crossings, \EvNAttrQp{} are attributed to catalogued
molecular lines, mostly carbon monoxide, \LnAttrSupp{} with independent support
(\S\ref{sec:mmdifferent}). Of the other \EvNUnattrQp{}, \CeTlObsCutWord{} is
atmospheric and \CeNTailCrossWord{} lie in the archival tail; the remaining
\CeObs{} compare with the \CeExp{} predicted from control positions beside the
star -- within what chance allows, and leaving room for a further \CeAddLim{}
(\S\ref{sec:trigunif}).

\item \emph{What the repeat observations show.} No crossing satisfies the
confirmation criterion, and the criterion is narrow: it reads the one cell a
discovery predicts, so it recovers only a carrier that returns to within half
a native channel -- as a spectral line does, and as the \RcNExcl{} tested,
accelerating at \ZtALosLo{} to \ZtALosHi\,m\,s$^{-2}$, do not unless a
periodic motion is caught at the matching phase. Where it applies it works:
\MkAttrRecur{} of the \MkAttrTested{} line-attributed crossings return.
Searched at every drift across each covering window, \WuNWinExc{} of the
\WuNWin{} windows reach the trigger against \WuExpWin{} expected
($P=\WuProbWin$; \WuExpWinCtrl{} at $P=\WuProbWinCtrl$ from control
positions), a marginal excess with \WuNWinExcTwo{} of it toward two objects;
the grid separates \TjNExcl{} of \TjNPair{} crossing pairs and leaves the rest
open. No transmitter is excluded by any of this. Confirmation is limited by
coverage -- \ZcNoRepSys{} systems have no repeat at a covered frequency -- and
by the criterion, not by depth (\S\ref{sec:mmdifferent}, \S\ref{sec:etacrv}).

\item \emph{What is not constrained, and what to do next.} The limits apply
only to a narrowband carrier outside the molecular-line mask, on while ALMA
looked and drifting by at most \DriftCeilLo\,nHz, \AccelSheikhPct{} per cent
of the recommended ceiling \citep{Sheikh2019}, and only to stars ALMA observed
(\S\ref{sec:excludes}). Other observers' choices of channels and repeats, not
collecting area, limited the search. A dedicated programme would channelise
finely, avoid the molecular lines and revisit each target at one tuning days
and years apart, beginning with the two objects that carry the excess:
HD~14055, the survey's most-observed star, and \BkAlmaName, whose share is
unexplained but not significant once its selection is allowed for (\S\ref{sec:twostars}, \S\ref{sec:future}).

\end{enumerate}

%% Five items: what is new, how deep, what was found, what the repeat
%% observations show, and what is not constrained.  Every number is cited from
%% the body; censusacct_r17 A4 and prose_v412 P9 read item 3.
